\documentclass[12pt]{amsart}

\usepackage[margin=1in]{geometry}
\usepackage{amsfonts}
\usepackage{amssymb}
\usepackage{hyperref}
\hypersetup{
    colorlinks=true,
    linkcolor=black,
    filecolor=green,
    urlcolor=cyan,
    citecolor=black,
    pdftitle={Teaching Statement},
}
\urlstyle{same}


\begin{document}

\begin{center}
  {\large\textsc{Teaching Statement}}\\[5pt]
  \small\textsc{Edward A. Varvak, University of Chicago}
\end{center}

My main goal in teaching is for students to leave a course able to explain why
the methods they learned work, and not only how to carry them out. As an undergrad, 
I was very frustrated by classes that prioritized memorization over an understanding 
of the material, and in my five years of teaching I have learned that I am not alone 
in this. In my own teaching, I try to show students what a formula is doing before
they use it. When I teach volumes of revolution, I draw the solid, sketch a typical
slice, and work out its area with the class, instead of handing students the washer
and shell formulas. Students then understand the integral as a sum of thin slices 
they can picture, as opposed to one more rule to memorize.


I have taught since 2022. At UIC, I was a teaching assistant for 
Finite Mathematics for Business (Math 160) in four semesters, and for
Calculus II (Math 181) in Spring 2026. This fall I am teaching Calculus II
(Math 15200) at the University of Chicago as a graduate student lecturer.
I teach one section of about thirty students, for which I write all lectures, 
problem sets, and exams myself. 

Most students in Finite Mathematics for Business take it to satisfy a requirement,
and many arrive unsure of their mathematical ability. Over several years of teaching it, 
I learned to temper my enthusiasm for more advanced applications of linear algebra
in favor of making sure my students were able to fully follow my explanations. 
Moving from TA to lecturer this fall has given me a much greater say in how my course 
is run. Where before I only needed to know the material for that day's section, 
I now plan every lecture myself, balancing the pace of the syllabus against 
my students' understanding.

My evaluators have frequently described my lecturing style as ``clear'' and ``engaging.''
In prior semesters at UIC, they often noted that 
I was an effective communicator, especially when engaging in one-on-one 
conversations with students. This fall, my evaluator praised 
the energy I brought to my 8:00 AM class, and remarked on his surprise at how engaged 
the students were. 


Beyond the classroom, I want to mentor students who are passionate about math.
My start in research was an undergraduate honors thesis with Olivier Martin, who 
was then a postdoc at Stony Brook. The experience pushed me to go into academia.
I would like to give undergraduates a similar start by mentoring reading projects 
in algebraic geometry and related topics. I am glad to teach the calculus sequence, 
linear algebra, and abstract algebra, and I would welcome the chance to teach an 
upper-level or graduate topics course in algebraic geometry or Hodge theory.

\end{document}
